\documentclass[unicode]{beamer}

% --- [設定] ---
\usepackage{luatexja}
\usepackage[orientation=portrait,size=a0,scale=1.4]{beamerposter}
\usepackage[most]{tcolorbox}
\usepackage{tikz}
\usetheme{default}

\setbeamertemplate{navigation symbols}{}
% ページ全体の基本余白を2cmに設定
\setbeamersize{text margin left=2cm, text margin right=2cm}
\addtobeamertemplate{block begin}{}{\setlength{\parskip}{0.5em}}
\addtobeamertemplate{block end}{}{\vspace{1em}} % ブロックが終わるたびに下に隙間を作る
\linespread{1.1}

% --- [フォントの設定：CloudLaTeX用] ---
\usepackage[haranoaji,deluxe]{luatexja-preset} % 原ノ味フォントを使用するプリセット
\renewcommand{\kanjifamilydefault}{\gtdefault}  % 日本語の標準をゴシック体にする
\renewcommand{\familydefault}{\sfdefault}      % 英数字の標準をサンセリフ（Arial風）にする

% --- [色の設定] ---
% ブロックの見出し（block title）の色をカスタマイズ
\setbeamercolor{block title}{fg=white, bg=cyan!70!black} % 文字を白、背景を濃い水色に
\setbeamercolor{block body}{bg=white}                   % 本文の背景を白に


% --- [数字を1桁なら全角にする命令] ---
\newcommand{\zenkakunum}[1]{%
  \ifnum#1<10%
    \ifcase#1 ０\or １\or ２\or ３\or ４\or ５\or ６\or ７\or ８\or ９\fi%
  \else%
    \number#1%
  \fi%
}

% 図の改良
\renewcommand{\figurename}{}
\renewcommand{\thefigure}{図\zenkakunum{\value{figure}}}

\makeatletter
\renewcommand{\p@figure}{}
\makeatother

\setbeamertemplate{caption}{\insertcaptionnumber ： \insertcaption}

% ブロック用のカウンター
\newcounter{blockcount}
% ブロックを自動ナンバリングで開始する命令
% 使い方: \begin{myblock}{研究背景} ．．． \end{myblock}
\newenvironment{myblock}[1]{%
  \refstepcounter{blockcount}%
  \begin{block}{\large \zenkakunum{\value{blockcount}}. #1}%
}{%
  \end{block}%
}


\begin{document}
\begin{frame}[t, fragile]

  % --- 1. タイトル部分 ---
  \vspace{1em} % 紙の一番上との間の隙間
  \nointerlineskip
  \begin{tcolorbox}[
      enhanced, sharp corners, boxrule=0pt,
      colback=cyan!20, 
      width=\textwidth,             % 本文の幅（左右余白を除いた幅）に合わせる
      % grow to left/right を消すことで、setbeamersize で設定した余白が活きます
      halign=center, height=7em
    ]
    {\Huge \bfseries 次元拡張} \par
    \vspace{1em}
    {\raggedleft 所属 \hspace{1cm} \par}
    {\raggedleft 著者 \hspace{1cm} \par}
  \end{tcolorbox}

  % --- 2. 中央の罫線（開始位置を調整） ---
  \begin{tikzpicture}[remember picture, overlay]
    \draw [gray!50, line width=3pt] 
      % yshift を -13cm くらいにすると、高さ10cmの箱の下に少し隙間が空きます
      ([yshift=-13cm]current page.north) -- ([yshift=2cm]current page.south);
  \end{tikzpicture}

  \vspace{0.1em} % タイトルボックスと本文の間の距離

  % --- 3. 2段組み ---
  \begin{columns}[t] 
    \begin{column}{.48\linewidth}
      \begin{myblock}{\large 研究動機}
研究動機 
\end{myblock}
      \begin{myblock}{\large 事前調査}
事前調査

・ド・グア

・三平方の拡張としての余弦定理とその意味付けについて
\end{myblock}
      \begin{myblock}{\large 仮説}
かせつ 
\end{myblock}
      \begin{myblock}{\large 方針}
放心
\end{myblock}
      \begin{myblock}{\large 結果}
けっか
\end{myblock}
    \end{column}

    \begin{column}{.02\linewidth}
      % 中央の線のスペース
    \end{column}

    \begin{column}{.48\linewidth}
      \begin{myblock}{\large 考察}
thinking...?
\end{myblock}
      \begin{myblock}{\large 結論}
結論です
\end{myblock}
      \begin{myblock}{\large 今後の展望}
4次元に拡張しよう、とかね
\end{myblock}
      \begin{myblock}{\large 参考文献}
さんこうにしたものをかいてください
\end{myblock}
    \end{column}
  \end{columns}

\end{frame}
\end{document}
